
\documentclass[11pt]{article}
\usepackage[margin=1in]{geometry}
\usepackage{amsmath}
\usepackage{graphicx}
\usepackage{hyperref}
\usepackage{booktabs}

\title{TNG Anisotropy Atlas: A Reproducible Pipeline for Satellite \\
Velocity Anisotropy in IllustrisTNG-100}
\author{[Your Name]}
\date{\today}

\begin{document}
\maketitle

\section{Goal}

Satellite galaxy velocity anisotropy,
$\beta(r) = 1 - \sigma_t^2(r) / (2\sigma_r^2(r))$,
characterizes the shape of satellite orbits within dark matter halos:
$\beta < 0$ indicates tangentially-biased orbits, $\beta > 0$ radially-biased
orbits, and $\beta = 0$ isotropy. This is well-studied in the literature
(e.g.\ Wojtak, Hansen \& Hjorth 2013, hereafter WHH13) and in standard
N-body simulations. The goal of this project is not to discover new
physics, but to build and validate a reproducible measurement pipeline
for $\beta(r)$ in IllustrisTNG-100, and to use it to produce a
population-level atlas across halo mass, environment, and galaxy
star-formation status. This pipeline and the resulting atlas are
intended as a validated baseline in standard cosmology, against which
future work testing alternative dynamical frameworks could be
compared -- no such comparison is attempted here, since IllustrisTNG
assumes standard gravity throughout.

\section{Data}

We use IllustrisTNG-100 (Nelson et al.\ 2019), snapshot 99 ($z=0$),
accessed via the public TNG API and the \texttt{illustris\_python}
package. Central host halos were selected in five mass bins
(Table~\ref{tab:massbins}), based on $\mathrm{Group\_M\_Crit200}$,
requiring at least 6 subhalos per host. Satellites were identified via
\texttt{SubhaloGrNr}, restricted to within $R_{200}$ of the host, with
positions and velocities computed relative to the host center under
periodic boundary conditions.

\begin{table}[h]
\centering
\begin{tabular}{lccc}
\toprule
Bin & $\log_{10}(M_{200}/10^{10}\,M_\odot/h)$ & $M_{200}$ range & $N_\mathrm{halos}$ \\
\midrule
0 & [1.0, 1.3) & $10^{11}$--$2\times10^{11}\,M_\odot/h$   & 3878 \\
1 & [1.3, 1.7) & $2\times10^{11}$--$5\times10^{11}$        & 3205 \\
2 & [1.7, 2.2) & $5\times10^{11}$--$1.6\times10^{12}$       & 1754 \\
3 & [2.2, 3.0) & $1.6\times10^{12}$--$10^{13}$              & 645  \\
4 & [3.0, 4.5) & $>10^{13}$ (cluster-scale)                 & 127  \\
\bottomrule
\end{tabular}
\caption{Host halo mass bins and sample sizes.}
\label{tab:massbins}
\end{table}

As a resolution check, we computed the halo mass function
$dn/d\log_{10}M$ for the full snapshot (Figure~\ref{fig:massfunc}).
The function turns over near $\log_{10}(M_{200}) \approx -1.6$,
consistent with the $\sim$50-particle resolution limit computed from
the simulation's dark matter particle mass
($m_\mathrm{DM} = 5.06\times10^6\,M_\odot/h$, sourced from the TNG
API). Our analysis sample begins at $\log_{10}(M_{200}) = 1.0$, roughly
2.6 dex above this limit, so halo-finding incompleteness does not
affect the mass bins used here.

\section{Methods}

\subsection{Kinematic decomposition}

For each satellite, position and velocity relative to the host center
are decomposed into radial and tangential components. Tangential
dispersion is computed as the second moment $\langle v_t^2 \rangle$,
not the variance of $v_t$: because tangential speed follows a
Rayleigh distribution, using the variance introduces a bias of
$+0.785(1-\beta)$, which we identified and corrected during toy-model
validation. Radial dispersion uses the variance (with $\mathrm{ddof}=1$)
around the bin mean, correcting for bulk flows. Errors are propagated
via the delta method.

\subsection{Validation on synthetic halos}

Before use on real data, the pipeline was validated on synthetic
halos with known input $\beta$ (Figure~\ref{fig:toy}). Across 30
random realizations each for $\beta_\mathrm{input} \in \{-0.5, 0, 0.5\}$,
recovered values matched inputs to within $\sim$1$\sigma$. A
radius-dependent test case, $\beta(r) = 0.5(r/r_\mathrm{max})^2$, was
also correctly recovered in shape.

\subsection{Stacking and bootstrap errors}

Satellites are stacked across all host halos within a mass bin, with
radius scaled by each host's $R_{200}$. Radial bins use equal-count
(quantile-based) edges. Uncertainties are estimated via bootstrap
resampling of \emph{host halos} (200 iterations), not individual
satellites, since satellites within a single halo are not independent
draws and satellite-level bootstrapping would understate the true
uncertainty.

\subsection{Environment and population splits}

Environment is measured as the number of neighboring massive halos
($\log_{10}(M_{200}) \geq 1.0$) within a 5~Mpc/h periodic search radius.
Star-forming/quenched classification uses a specific star formation
rate threshold of $10^{-11}\,\mathrm{yr}^{-1}$ (Franx et al.\ 2008).

\section{Validation}

Figure~\ref{fig:massfunc} shows the halo mass function resolution
check described above. Figure~\ref{fig:quenchradial} shows quenched
fraction as a function of clustercentric radius in the two
galaxy-resolved bins (Bin 3, $N=302$ halos; Bin 4/clusters, $N=127$
halos), declining monotonically from 65\% to 31\% (Bin 3) and 94\% to
70\% (Bin 4) from center to outskirts. This trend is well-established
in the galaxy evolution literature (ram-pressure stripping and tidal
quenching are most efficient near cluster cores); we do not claim it
as novel, but present it as an independent, observable-space
validation that the pipeline reproduces known physics.

\section{Results}

\subsection{Mass trend}

The core result (Figure~\ref{fig:betamass}) is a monotonic trend from
tangentially-biased orbits at group scale to radially-biased orbits at
cluster scale: Bin 0, $\beta \approx -0.32$ to $-0.03$; Bin 1,
$\approx -0.18$ to $-0.05$; Bin 2, $\approx -0.18$ to $-0.03$; Bin 3,
$\approx -0.06$ to $0.02$ (roughly isotropic); Bin 4 (clusters),
$\approx +0.07$ to $+0.24$ (radially biased). This is consistent with
low-mass halos being dynamically relaxed with tidally-circularized
orbits, and cluster halos actively assembling via radial filamentary
infall -- matching the qualitative expectation from WHH13 and standard
N-body halo anisotropy profiles ($\beta \approx 0$ at the center,
rising to 0.4--0.6 in the outskirts).

\subsection{Environment}

Per-halo Spearman correlations ($N=2580$ individual halos with valid
$\beta$) give $\rho(\beta, \mathrm{mass}) = 0.139$
($p = 1.55\times10^{-12}$) and, controlling for mass via partial
Spearman correlation, $\rho(\beta, \mathrm{environment}\,|\,\mathrm{mass})
= -0.099$ ($p = 4.19\times10^{-7}$): environment has a statistically
significant effect on $\beta$ independent of mass.

A median environment split within each mass bin shows a sign reversal:
at group/intermediate mass, denser environments correspond to
\emph{more} tangential bias ($\Delta\beta$ from $-0.073$ to $-0.084$
across Bins 0--3), while at cluster mass, denser environments
correspond to \emph{more} radial bias ($\Delta\beta = +0.048$ in Bin
4). We interpret this as increased tidal stirring at group scale
versus enhanced filamentary infall feeding actively-growing clusters,
but note Bin 4's environment split has only 62--65 halos per half; we
present this as a suggestive, physically-motivated finding warranting
further investigation with larger-volume simulations, not a confirmed
result.

\subsection{Resolution limit and galaxy-only cross-check}

The vast majority of subhalos in our catalog are dark-matter-only,
with zero resolved stellar mass (95.2\%--99.4\% depending on mass
bin) -- a known TNG100 resolution limit for dwarf galaxies, not a
pipeline bug. Applying a resolved-galaxy stellar mass floor of
$10^{8.5}\,M_\odot/h$, only Bins 2--4 retain usable samples (20, 302,
and 127 halos respectively; Bins 0--1 have zero halos with $\geq 3$
resolved satellites). Galaxy-only satellites are systematically more
radially biased than the full subhalo population (e.g.\ Bin 3:
$-0.019 \to +0.144$; Bin 4: $+0.162 \to +0.278$), consistent with
stronger dynamical friction acting on more massive, luminous
satellites.

\section{Limitations}

\begin{itemize}
\item The Bin 4 (cluster) sample is small ($N=127$ halos), and the
  environment-split sign reversal in this bin should be treated as
  suggestive rather than confirmed.
\item Bin-level Spearman correlations (across only 5 mass bins) are
  statistically weak; the per-halo correlation ($N=2580$) is the
  defensible statistic and is what we report as significant.
\item TNG100's mass resolution prevents resolving dwarf galaxies in
  group-mass hosts, restricting the galaxy-only cross-check to Bins
  2--4.
\item All results assume standard gravity, as encoded in the
  IllustrisTNG simulation physics.
\end{itemize}

\section{Baseline for alternative dynamics}

This pipeline and atlas are presented purely as a validated baseline
within standard $\Lambda$CDM cosmology and standard gravity, as
implemented in IllustrisTNG. No claims are made, implied, or intended
regarding modified gravity, modified inertia, or any alternative
dynamical framework. Future work seeking to test such frameworks
against satellite kinematics would require either a simulation run
under the alternative physics, or an independent observational
comparison; this atlas is not designed for and should not be
interpreted as evidence for or against any such framework.

\end{document}
